\chapter{Binary search}\label{ch:bsearch}
\index{binary search}

Chapter~3 spent effort rearranging an input until it looked the way a
greedy needed it. This chapter rearranges nothing at all: it takes a
question whose answer is \emph{one number}, notices that the number's
candidates split cleanly into ``too small'' and ``big enough'', and finds
the boundary in about sixty probes --- whatever the range. Most of the
problems a Silver contestant calls hard are this idea wearing a costume:
the answer to \emph{what is the smallest makespan} or \emph{what is the
largest safe gap} is easier to \emph{verify} than to construct, and
verification is enough. After this chapter you will be able to write the
one template that never loops, spot a monotone predicate inside a word
problem, and turn ``minimise the maximum'' into a routine search over the
answer.

\section{The one correct template}\label{sec:bs-template}
\index{binary search!template}

Binary search has a terrible reputation for off-by-one bugs, and every
version of it deserves the reputation except this one. The trick is to
search \emph{predicates}, not elements: pick a question $P(x)$ that is
false for a while and then true forever, and find the first true.
On a sorted array the question ``$a[i] \ge t$'' is exactly that shape.

\begin{lstlisting}[language=C++,caption={\texttt{lib/bsearch.h}: first
true in ${[lo,hi]}$, given a predicate monotone on the whole
interval.},label={lst:bsearch}]
template <class F> ll first_true(ll lo, ll hi, F pred) {
    while (lo < hi) {
        ll mid = lo + (hi - lo) / 2;
        if (pred(mid)) hi = mid; else lo = mid + 1;
    }
    return lo;
}
\end{lstlisting}

Three invariants carry the whole proof of correctness, and you should be
able to recite them. \textbf{(a)} Every position in $[lo,hi)$ is still a
candidate for the answer. \textbf{(b)} Everything left of $lo$ is known
false, everything at or right of $hi$ is known true.
\textbf{(c)} $hi - lo$ strictly shrinks, so the loop ends with
$lo = hi$: one candidate, the first true.
Two details do the real work: $\texttt{mid}$ rounds \emph{down}
(\texttt{lo + (hi - lo)/2}, which also dodges overflow), and the update
is \texttt{hi = mid} on true versus \texttt{lo = mid + 1} on false ---
the false side can skip \texttt{mid} entirely, the true side cannot.
\index{binary search!invariants}

The standard library sells the same idea as \texttt{lower\_bound}, which
returns an iterator to the first element $\ge t$; upper\_bound returns
the first $> t$. Both are the first-true shape.

\begin{lstlisting}[language=C++,caption={\texttt{bsearch\_template.cpp}:
the template and \texttt{lower\_bound} agree.},label={lst:bstemp}]
\end{lstlisting}
\lstinputlisting{code/bsearch_template.cpp}

Real output on $a = \{2,4,7,7,11\}$, $t = 7$:
\begin{lstlisting}[numbers=none]
first >= 7 at index 2 (a[2]=7)
first >  7 at index 4 (a[4]=11)
lower_bound agrees: 2
\end{lstlisting}

\begin{figure}[ht]\centering
% fig15-frontier.tex -- Figure 4.1: predicate frontier F...F T...T for the factory check at (3 7 / 3 2 5).
% fstyle.tex -- shared TikZ styles for the figures of this book (palette from lifeos.sty).
\tikzset{%
  cell/.style={draw=ink, minimum width=7mm, minimum height=6.5mm, inner sep=0pt, font=\footnotesize, fill=white},
  cellh/.style={cell, fill=accentsoft, draw=accent},
  cella/.style={cell, fill=amberbg, draw=amber},
  cellg/.style={cell, fill=goodbg, draw=good},
  cellx/.style={cell, fill=graybg, draw=slate, text=slate},
  vtx/.style={circle, draw=ink, minimum size=6.5mm, inner sep=0pt, font=\footnotesize, fill=white},
  vtxh/.style={vtx, draw=accent, fill=accentsoft},
  vtxa/.style={vtx, draw=amber, fill=amberbg},
  vtxg/.style={vtx, draw=good, fill=goodbg},
  vtxd/.style={vtx, draw=slate, fill=graybg, text=slate},
  lab/.style={font=\footnotesize, text=slate},
  labA/.style={font=\footnotesize, text=accent},
  labM/.style={font=\footnotesize, text=amber},
  idx/.style={font=\tiny, text=slate},
  arr/.style={-{stealth}, thick, ink},
  arrA/.style={-{stealth}, thick, accent},
  arrM/.style={-{stealth}, thick, amber},
  arrG/.style={-{stealth}, thick, good},
  edge/.style={thick, ink},
  edgeA/.style={thick, accent},
  edgeM/.style={thick, amber},
}%
% \cellrow{prefix}{x0}{y}{style}{v0,v1,...}: a row of cells named prefix0, prefix1, ... (pitch 0.75cm)
\providecommand{\cellrow}[5]{\foreach \v [count=\i from 0] in {#5} {\node[#4] (#1\i) at (#2+\i*0.75,#3) {\v};}}%
% \cellidx{prefix}{n}{start}: index labels under cells prefix0 .. prefix(n-1)
\providecommand{\cellidx}[3]{\foreach \i in {0,...,\numexpr#2-1\relax} {\node[idx, below=0pt of #1\i] {\the\numexpr\i+#3\relax};}}%
% \fitfig{<tikzpicture>}: shrink the picture to \linewidth only when it would overflow.
\ifdefined\fitbox\else\newsavebox{\fitbox}\fi%
\providecommand{\fitfig}[1]{\sbox{\fitbox}{#1}\ifdim\wd\fitbox>\linewidth\resizebox{\linewidth}{!}{\usebox{\fitbox}}\else\usebox{\fitbox}\fi}%

\begin{tikzpicture}[font=\footnotesize]
\cellrow{c}{0}{0}{cella, minimum width=0.95cm}{1,3,5,7}
\foreach \i/\v in {4/9,5/11,6/13,7/15} \node[cellh, minimum width=0.95cm] (c\i) at (\i*0.75,0) {\v};
\foreach \i in {0,...,3} \node[font=\scriptsize, text=amber, below=1pt of c\i] {F};
\foreach \i in {4,...,7} \node[font=\scriptsize, text=accent, below=1pt of c\i] {T};
\draw[decorate,decoration={brace,raise=3pt},amber] (c0.north west) -- (c3.north east) node[midway,above=5pt,labM]{reject side (F)};
\draw[decorate,decoration={brace,raise=3pt},accent] (c4.north west) -- (c7.north east) node[midway,above=5pt,labA]{accept side (T)};
\draw[arrA] ($(c3.south east)!0.5!(c4.south west)+(0,-0.65)$) -- ++(0,0.58);
\node[labA, align=center, below] at ($(c3.south east)!0.5!(c4.south west)+(0,-0.65)$) {the answer is the \emph{first} T\\[-1pt]the frontier};
\end{tikzpicture}

\caption{The predicate frontier for $P(T)=$ ``can we finish by time
$T$'' on the Factory Machines instance. Everything the search does is
walk the boundary between amber and blue --- and it commits to the
\emph{first} T, never any T to its right.}
\label{fig:frontier}
\end{figure}

\begin{figure}[ht]\centering
\begin{tikzpicture}[font=\footnotesize]
% F4.3 input space vs answer space.
\node[draw=slate,fill=graybg,align=left,text width=4.6cm] at (0,0.6)
  {\textbf{input space}: the array the problem hands you ---
   $n$ up to $2\cdot10^5$, read once};
\node[draw=accent,fill=accentsoft,align=left,text width=4.6cm] at (0,-1.2)
  {\textbf{answer space}: the value the problem asks for ---
   a number you \emph{choose}, so its shape is yours};
\draw[{stealth}-,thick,ink] (-2.6,-2.6) -- (2.6,-2.6);
\foreach \t/\lbl in {-2/1, -0.7/, 0.6/, 2/{$\sim10^{18}$}}
  \draw (\t,-2.7) -- (\t,-2.5);
\node[below,font=\tiny] at (-2,-2.75) {1};
\node[below,font=\tiny] at (2,-2.75) {$\sim 10^{18}$};
\node[accent] at (0.6,-2.25) {$T^{\ast}$};
\draw[{stealth}-,accent,thick] (0.6,-2.5) -- ++(0,-0.7)
  node[below,align=center,text=accent,text width=5.4cm]
  {$\log_2$ of this whole line is about 60 probes\\
   --- the answer space shrinks like the array did};
\end{tikzpicture}

\caption{Two different spaces. The input array is data you must read;
the answer space is a line you get to choose the probing on ---
and $\log_2$ of $10^{18}$ is only about 60.}
\label{fig:answerspace}
\end{figure}

\begin{figure}[ht]\centering
% fig96-bswalk.tex -- Figure 4.x: first_true walk on a monotone predicate over positions 1..8 (first true at 6).
% fstyle.tex -- shared TikZ styles for the figures of this book (palette from lifeos.sty).
\tikzset{%
  cell/.style={draw=ink, minimum width=7mm, minimum height=6.5mm, inner sep=0pt, font=\footnotesize, fill=white},
  cellh/.style={cell, fill=accentsoft, draw=accent},
  cella/.style={cell, fill=amberbg, draw=amber},
  cellg/.style={cell, fill=goodbg, draw=good},
  cellx/.style={cell, fill=graybg, draw=slate, text=slate},
  vtx/.style={circle, draw=ink, minimum size=6.5mm, inner sep=0pt, font=\footnotesize, fill=white},
  vtxh/.style={vtx, draw=accent, fill=accentsoft},
  vtxa/.style={vtx, draw=amber, fill=amberbg},
  vtxg/.style={vtx, draw=good, fill=goodbg},
  vtxd/.style={vtx, draw=slate, fill=graybg, text=slate},
  lab/.style={font=\footnotesize, text=slate},
  labA/.style={font=\footnotesize, text=accent},
  labM/.style={font=\footnotesize, text=amber},
  idx/.style={font=\tiny, text=slate},
  arr/.style={-{stealth}, thick, ink},
  arrA/.style={-{stealth}, thick, accent},
  arrM/.style={-{stealth}, thick, amber},
  arrG/.style={-{stealth}, thick, good},
  edge/.style={thick, ink},
  edgeA/.style={thick, accent},
  edgeM/.style={thick, amber},
}%
% \cellrow{prefix}{x0}{y}{style}{v0,v1,...}: a row of cells named prefix0, prefix1, ... (pitch 0.75cm)
\providecommand{\cellrow}[5]{\foreach \v [count=\i from 0] in {#5} {\node[#4] (#1\i) at (#2+\i*0.75,#3) {\v};}}%
% \cellidx{prefix}{n}{start}: index labels under cells prefix0 .. prefix(n-1)
\providecommand{\cellidx}[3]{\foreach \i in {0,...,\numexpr#2-1\relax} {\node[idx, below=0pt of #1\i] {\the\numexpr\i+#3\relax};}}%
% \fitfig{<tikzpicture>}: shrink the picture to \linewidth only when it would overflow.
\ifdefined\fitbox\else\newsavebox{\fitbox}\fi%
\providecommand{\fitfig}[1]{\sbox{\fitbox}{#1}\ifdim\wd\fitbox>\linewidth\resizebox{\linewidth}{!}{\usebox{\fitbox}}\else\usebox{\fitbox}\fi}%

\begin{tikzpicture}[font=\footnotesize]
\node[lab, anchor=east] at (-0.55,0) {$P(x)$};
\foreach \i/\v in {1/F,2/F,3/F,4/F,5/F,6/T,7/T,8/T}
  \node[\ifnum\i<6 cella\else cellh\fi, minimum width=8mm, minimum height=7mm] (c\i) at (\i*1.0,0) {\v};
\foreach \i in {1,...,8} \node[idx, below] at (c\i.south) {\i};
\node[labA, text width=2.4cm, align=center, anchor=north] (p1) at (3.2,-1.2) {probe 1: $\mathit{mid}=4$ is F\\ $\Rightarrow$ throw the left half};
\node[labA, text width=2.4cm, align=center, anchor=north] (p2) at (7.2,-1.2) {probe 2: $\mathit{mid}=7$ is T\\ $\Rightarrow$ keep its left side};
\node[labA, text width=6.4cm, align=center, anchor=north] (p3) at (5.0,-3.1) {probe 3: $\mathit{mid}=5$ is F;\\ probe 4: $\mathit{mid}=6$ is T: \emph{first true}};
\draw[arrA] ($(p1.north)+(0.8,0)$) -- (c4.south); \draw[arrA] (p2.north) -- (c7.south); \draw[arrA] (5,-3.1) -- (5,-0.85);
\draw[ink, very thick] (c5.north east) -- ++(0,0.4) node[above, labA] {boundary};
\end{tikzpicture}

\caption{A first\_true walk on the predicate row of
Figure~\ref{fig:frontier}: each probe lands on a midpoint, the
surviving half keeps the boundary inside it, and four probes suffice on
eight cells --- $\lceil\log_2 8\rceil + 1$. Watch the invariant, not
the values: the true/false boundary never leaves the live interval.}
\label{fig:bswalk}
\end{figure}


\begin{figure}[htbp]\centering
\begin{tikzpicture}[font=\scriptsize]
\foreach \x/\v in {0/F,1/F,2/F,3/F,4/F,5/T,6/T,7/T,8/T}{
  \ifx\v F\draw[very thick,fill=amber!25] (\x*0.72,0) rectangle (\x*0.72+0.7,0.6);
  \else\draw[very thick,fill=accent!20] (\x*0.72,0) rectangle (\x*0.72+0.7,0.6);\fi
  \node at (\x*0.72+0.35,0.3) {\v};}
\draw[->,very thick] (5.4,1.2) -- (3.95,1.2) node[above,pos=0]{first true};
\node[align=left] at (3.2,-0.9) {monotone predicate:\\the boundary is the answer};
\end{tikzpicture}

\caption{The predicate picture: $F$ everywhere left, $T$ right --- binary search finds the boundary.}
\label{fig:predicate}
\end{figure}

\section{Monotone predicates: searching the answer}\label{sec:bs-answer}
\index{binary search!on the answer}\index{monotone predicate}

Figure~\ref{fig:frontier} is the whole idea; Figure~\ref{fig:bswalk} walks it once, probe by probe. When the predicate is
monotone --- false up to some boundary, true ever after (or the reverse)
--- the search runs on the \emph{answer space}, the interval of every
value the answer could take, not on the input array. Three questions
design the search and nothing else matters:

\begin{enumerate}
\item \textbf{What is the answer's range?} $[lo, hi]$ must provably
contain it --- the fastest machine's time times $t$ is a safe top.
\item \textbf{What is the check?} $P(x)$ must be computable in far less
time than constructing the answer.
\item \textbf{Is it monotone?} If $P$ is a sawtooth, the search finds
\emph{a} boundary, not \emph{the} answer (Section~\ref{sec:bs-limits}).
\end{enumerate}

\begin{example}[CSES \emph{Factory Machines}, min makespan]\label{ex:factory}
\index{Factory Machines}
$n$ machines make products, machine $i$ taking $k_i$ seconds each.
Shortest time for $t$ products total ($n \le 2\cdot10^5$, $t,k_i \le
10^9$)?
\par\smallskip
\textbf{Strategy.} The answer is a time $T$; the answer space is
$[1, 10^{18}]$. The predicate ``can we finish by $T$'' is monotone ---
more time never hurts --- and cheap: machine $i$ produces
$\lfloor T/k_i \rfloor$ units, so $\sum_i \lfloor T/k_i \rfloor \ge t$
decides it in $O(n)$. Search over ~60 probes $\times$ one sum each.
\par\smallskip
\begin{solution}
\begin{lstlisting}[language=C++,caption={\texttt{factory\_machines.cpp}:
the check plus the driver.},label={lst:factory}]
\end{lstlisting}
\lstinputlisting{code/factory_machines.cpp}
\end{solution}
Real output (CSES sample: $3$ machines $\{3,2,5\}$, $t=7$):
\begin{lstlisting}[numbers=none]
check(7)=0 check(8)=1
8
\end{lstlisting}
The check itself shows the frontier: 7 fails, 8 passes, and 8 is the
answer. One overflow habit to steal: the accumulator stops at $t$
instead of summing $10^9$ terms it does not need.
\end{example}

\begin{figure}[ht]\centering
\begin{tikzpicture}[font=\footnotesize,x=0.62cm,y=0.62cm]
% F4.4 factory check as a machine-completion timeline (T=8, rates 3,2,5).
\draw[{stealth}-,ink] (0,0) -- (9.2,0) node[below left]{time};
\foreach \t in {0,...,8} \draw[ink] (\t,-0.08) -- (\t,0.08)
  node[below=2pt,font=\tiny]{\t};
\foreach \y/\rate/\lbl in {1.0/3/$k_1{=}3$, 2.0/2/$k_2{=}2$, 3.0/5/$k_3{=}5$} {
  \node[anchor=east] at (-0.2,\y) {\lbl};
  \draw[slate] (0,\y) -- (8.4,\y);
}
\foreach \x in {3,6} \fill[accent] (\x,1.0) circle (2.2pt);
\foreach \x in {2,4,6,8} \fill[accent] (\x,2.0) circle (2.2pt);
\foreach \x in {5} \fill[accent] (\x,3.0) circle (2.2pt);
\draw[amber,thick,dashed] (8,0.15) -- (8,3.5)
  node[above,text=amber]{$T{=}8$: $2{+}4{+}1=7$};
\end{tikzpicture}

\caption{The check for Example~\ref{ex:factory} as a completion
timeline at $T=8$: machine $k_1{=}3$ finishes at $3,6$ (two products),
$k_2{=}2$ at $2,4,6,8$ (four), $k_3{=}5$ at $5$ (one). Counting ticks
left of the dashed line \emph{is} the predicate --- nothing was
scheduled.}
\label{fig:timeline}
\end{figure}

\begin{example}[CSES \emph{Array Division}, minimise the maximum
sum]\label{ex:arrdiv}\index{Array Division}
Split an array of $n$ positive values into $k$ contiguous groups
minimising the largest group sum ($n \le 2\cdot10^5$).
\par\smallskip
\textbf{Strategy.} Answer space: $[\max a_i,\ \sum a_i]$. Predicate
``does cap $x$ per group suffice'' is monotone --- a bigger cap can only
help --- and checkable greedily: accumulate until the next element would
overflow the cap, cut, count groups.
\par\smallskip
\begin{solution}
\lstinputlisting{code/array_division.cpp}
Real output ($n{=}5$, $k{=}3$, $\{2,4,7,3,5\}$):
\begin{lstlisting}[numbers=none]
check(7)=0 check(8)=1
8
\end{lstlisting}
The optimal split is $[2,4],[7],[3,5]$ --- but notice the program never
built it. The greedy's correctness rests on a minimality claim: each
greedy group extends as far right as the cap allows, so no feasible
division uses fewer groups. If the greedy count $g$ is at most $k$,
merging adjacent groups --- always legal, each merge just concatenates
two pieces --- reduces the count to exactly $k$; if $g$ exceeds $k$, no
division exists at all. Asking only ``$g \le k$'' is therefore sound.
\end{solution}
\end{example}

\begin{example}[USACO-S \emph{Social Distancing}, maximise the
minimum]\label{ex:social}\index{Social Distancing}
Free stalls at sorted positions $p_1 < \dots < p_n$; place $k$ cows to
maximise the smallest gap between neighbours.
\par\smallskip
\textbf{Strategy.} The min-max pattern inverted: now the answer is the
\emph{largest} feasible gap $d$. Predicate ``can $k$ cows sit with gaps
$\ge d$'' is monotone the other way --- feasible for small $d$, infeasible
past the frontier. The check is the same greedy as ever: place the next
cow at the first stall at least $d$ past the last.
\par\smallskip
\begin{solution}
\lstinputlisting{code/social_distancing.cpp}
Real output (stalls $\{1,2,5,8,9\}$, $k=3$):
\begin{lstlisting}[numbers=none]
fits(4)=1 fits(5)=0
4
\end{lstlisting}
Two habits worth copying: maximisation becomes first-true on the
\emph{negated} predicate (search for the first $d$ that fails, answer is
one less), and the candidate range runs to the span of the stalls ---
answers beyond it are impossible by construction.
\end{solution}
\end{example}

\begin{example}[CSES \emph{Multiplication Table}, a counting
predicate]\label{ex:multtable}\index{Multiplication Table}
The $n \times n$ multiplication table ($n$ odd, $n < 10^6$): find the
median of its $n^2$ products without building them.
\par\smallskip
\textbf{Strategy.} Nothing is sorted and nothing can be --- $n^2$ is too
many elements to hold. The answer space is the product range
$[1, n^2]$, and the predicate is a \emph{count}: ``at least
$\lceil n^2/2 \rceil$ products are $\le x$''. Row $i$ contributes
$\min(n, \lfloor x/i \rfloor)$, so counting is $O(n)$ and monotone ---
larger $x$ can only add products.
\par\smallskip
\begin{solution}
\lstinputlisting{code/mult_table.cpp}
Real output:
\begin{lstlisting}[numbers=none]
3
\end{lstlisting}
$n{=}3$: the sorted products are $\{1,2,2,3,3,4,6,6,9\}$ and the middle
is $3$. The pattern to keep: \emph{rank queries} (``how many are $\le
x$'') are monotone almost by definition, so any ``find the $k$-th''
question is candidate answer-space search.
\end{solution}
\end{example}

\begin{example}[CSES \emph{Concert Tickets}: the STL's own first\_true]
\label{ex:concert}
Tickets cost $h_1,\dots,h_n$; each of $m$ customers in turn names a
maximum price $t_i$ and takes the most expensive ticket costing at
most $t_i$, or none ($n, m \le 2\cdot10^5$).
\par\smallskip
\textbf{Strategy.} Sort the tickets and binary-search each $t_i$ ---
that is the template verbatim, but the STL already carries it:
a \texttt{multiset} keeps prices sorted, and
\texttt{upper\_bound($t$)} returns \emph{the first element exceeding}
$t$ --- the same first\_true boundary as
\S\ref{sec:bs-template}, with ``too expensive'' as the predicate.
Step back one iterator for the answer. The custom walk still matters:
here the library does it for you, and knowing that is the difference
between a three-line solution and a forty-line one.
\par\smallskip
\begin{solution}
\texttt{code/concert\_tickets.cpp}:
\begin{lstlisting}[numbers=none]
$ ./concert_tickets < tickets_sample.txt
3
8
-1
\end{lstlisting}
Customer three wants $\le 3$ but only $\{5,7,5\}$ remain ---
\texttt{upper\_bound} lands on \texttt{begin()} and we print $-1$.
Each query is $O(\log n)$; sorted input plus an ordered container is
the whole trick.
\end{solution}
\end{example}

\begin{exercise}[easy]\label{exr:boundary}
The template keeps ``left of \texttt{lo} known false, at-or-right-of
\texttt{hi} known true''. For each of these predicate answers at
\texttt{mid}, state the correct update: (a) $P(\texttt{mid})$ false,
(b) $P(\texttt{mid})$ true, (c) the invariant when the loop exits.
\end{exercise}

\begin{exercise}[medium]\label{exr:monotone}
Which of these predicates is monotone, and on what range?
(a) ``$x$ machines can finish by time $T$'' as a function of $x$;
(b) ``$x^2$ is prime''; (c) ``the greedy's group count is $\le k$'';
(d) ``$a[i] \ge t$'' on an \emph{unsorted} array.
\end{exercise}

\section{Real-valued search}\label{sec:bs-real}
\index{binary search!real-valued}

Over a continuous interval the ``first true'' trick never terminates ---
there is no first true real number. The fix is to stop after a fixed
number of iterations: each probe halves the bracket, so 60 rounds shrink
$[0,2]$ by $2^{-60} \approx 10^{-18}$, well past double precision.
Never write \texttt{while (hi - lo > eps)}: it compares the bracket to a
tolerance that ignores the function's steepness, and \texttt{mid*mid}
near 1e20 can overflow silently. Fixed iterations are cheaper to reason
about and just as accurate.

\begin{lstlisting}[language=C++,caption={\texttt{real\_bsearch.cpp}:
60 iterations to $\sqrt 2$ --- the bracket, not the midpoint, is the
deliverable.},label={lst:real}]
\end{lstlisting}
\lstinputlisting{code/real_bsearch.cpp}
Real output:
\begin{lstlisting}[numbers=none]
it=10 bracket=[1.4140625, 1.416015625]
it=20 bracket=[1.4142131805419922, 1.414215087890625]
it=60 bracket=[1.4142135623730949, 1.4142135623730951]
\end{lstlisting}

Note which side moved: the predicate is $x^2 < 2$, true below the root
--- so \texttt{lo} chases the boundary from below and \texttt{hi} from
above. Reading off which endpoint is safe to report is the whole
correctness story.


\begin{figure}[htbp]\centering
\begin{tikzpicture}[font=\scriptsize]
\foreach \x/\v in {0/3,1/5,2/5,3/8,4/9,5/12}{
  \draw[very thick,fill=accent!12] (\x*0.85,0) rectangle (\x*0.85+0.8,0.6);
  \node at (\x*0.85+0.4,0.3) {\v};}
\draw[->,very thick,amber] (2.5,1.3) -- (3.4,0.75);
\node[align=left] at (2.5,-0.9) {lower\_bound(8) $\to$ first $\ge 8$:\\skips the two 5s, lands on 8};
\end{tikzpicture}

\caption{\texttt{lower\_bound} on a sorted multiset: the first element $\ge$ the key.}
\label{fig:multiset}
\end{figure}

\section{Where the pattern ends}\label{sec:bs-limits}
\index{binary search!limitations}

The template never fails; the \emph{predicate} fails. A non-monotone
$P$ still returns a boundary, but it is the boundary of wherever
\texttt{mid} happened to land --- the answer space may hold several
frontiers, and you find whichever the search path crossed first. Two
honest warnings: (a) monotonicity is a \emph{proof obligation}, not a
property you get for free from writing \texttt{check()}; convince
yourself the frontier is unique before you search it.
(b) A monotone-sounding English description can still hide a sawtooth:
``$x$ workers finish in time'' is monotone; ``configuration $x$ is
valid'' need not be. If the predicate wiggles, you need structure ---
convexity, a transform --- not optimism.

\section{Chapter summary}
\begin{table}[ht]\centering
{\small
\begin{tabularx}{\linewidth}{@{}lYYYY@{}}
\toprule
\tablehead\thead{Problem} & \thead{Answer space} & \thead{Predicate} &
\thead{Check cost} & \thead{WE}\\
\midrule
Factory Machines & $[1,\,10^{18}]$ time & $\sum \lfloor T/k_i\rfloor\ge t$ &
$O(n)$ & \ref{ex:factory}\\
Array Division & $[\max,\sum]$ cap & greedy groups $\le k$ & $O(n)$ &
\ref{ex:arrdiv}\\
Social Distancing & $[0,\text{span}]$ gap & greedy placing $\ge k$ &
$O(n)$ & \ref{ex:social}\\
Multiplication Table & $[1,n^2]$ value & count $\le x \ge$ median &
$O(n)$ & \ref{ex:multtable}\\
$\sqrt 2$ on $\mathbb R$ & $[0,2]$ real & $x^2 < 2$ & $O(1)$ &
\S\ref{sec:bs-real}\\
\bottomrule
\end{tabularx}}
\caption{One template, five costumes. In every row the hard work was
\emph{writing the check} --- the search was already solved.}
\label{tab:bsearch}
\end{table}

\section{Common mistakes}
\index{common mistakes}
\begin{itemize}
\item \textbf{Mixing the updates.} \texttt{hi = mid - 1} on true throws
away the only candidate you know is true; \texttt{lo = mid} on false
leaves \texttt{mid} in play and loops forever. Memorise the pair:
\texttt{true} keeps \texttt{mid}, \texttt{false} skips it.
\item \textbf{Searching the wrong direction.} The template finds the
first \emph{true}. For a last-true or first-false problem, negate the
predicate or subtract one --- do not improvise a mirrored loop.
\item \textbf{Predicate that is not monotone.} The search still returns
something; it is just not the answer you wanted. Prove the frontier is
unique first.
\item \textbf{Overflow inside the check.} $\sum T/k_i$ with
$T = 10^{18}$ and $n = 2\cdot10^5$ overflows a 64-bit sum several times
over; cap the accumulator at the target, as Listing~\ref{lst:factory}
does.
\item \textbf{\texttt{while (hi-lo > eps)} on reals.} Epsilons couple
the stop rule to the function's slope and invite infinite loops; fix
the iteration count instead.
\end{itemize}

\section{Additional exercises}
\begin{exercise}[easy]
CSES \emph{Sum of Three Values} --- reduce to two-sum via sorting.
\end{exercise}
\begin{exercise}[easy]
USACO-S \emph{Cow Dance Show} --- feasibility as a simulation per
candidate size.
\end{exercise}
\begin{exercise}[easy]
USACO-S \emph{Convention} --- monotone coverage check.
\end{exercise}
\begin{exercise}[medium]
CSES \emph{Reading Books} --- the answer space is a waiting time.
\end{exercise}
\begin{exercise}[medium]
CSES \emph{Movie Festival II} --- multiset membership per step; where
does the monotonicity live?
\end{exercise}
\begin{exercise}[medium]
USACO-S \emph{Loan Repayment} --- search the repayment term.
\end{exercise}
\begin{exercise}[medium]
USACO-S \emph{Bakery} --- feasibility linear program per cost bound.
\end{exercise}
\begin{exercise}[medium]
USACO-S \emph{Angry Cows} --- the classic blast-radius search.
\end{exercise}
\begin{exercise}[medium]
CF \emph{Guess the K-th Zero} --- an \emph{interactive} first-true.
\end{exercise}
\begin{exercise}[easy] Hand-walk the first\_true template on predicate \texttt{F F F T T} and list every probe. \end{exercise}
\begin{exercise}[medium] USACO-S \emph{Angry Cows} --- search the blast radius directly. \end{exercise}
\begin{exercise}[hard]
CSES \emph{Maximum Subarray Sum II} --- combine with a sliding window.
\end{exercise}
\begin{exercise}[hard]
CF \emph{The Meeting Place Cannot Be Changed} --- real-valued answer.
\end{exercise}
\begin{exercise}[hard]
Baltic OI 2012 \emph{Mobile} --- heavy feasibility per candidate weight.
\end{exercise}

\chapternotes
Binary search on the answer is the flagship technique of
Sannemo's chapter on it~\cite{sannemo2025} and Halim's greedy/search
sections~\cite{halimcp4}; Laaksonen's guide covers the template's
pitfalls~\cite{laaksonen2024}. The four archive problems are official
CSES/USACO tasks~\cite{cses,usacoguide}; standard references on the
idiom are the C++ reference's \texttt{lower\_bound}
page~\cite{cppreference} and the cp-algorithms binary-search note.
\begin{exercise}[hard] CSES \emph{Sliding Median} --- order statistics inside the window. \end{exercise}

\section*{Further reading}
\begin{itemize}
\item USACO Guide, ``Binary Search'' (intermediate) --- same first\_true
template, more answer-search problems at the Silver level. \newline\url{https://usaco.guide/silver/binary-search}
\item Laaksonen, \emph{Competitive Programmer's Handbook}, ch.~4
``Binary search'' (beginner) ---
the array half of the technique plus STL lower\_bound idiom. \newline\url{https://cses.fi/book/book.pdf}
\item cp-algorithms, ``Binary search'' (reference) ---
implementation variants incl.\ the real-valued search of
\S\ref{sec:bs-real}. \newline\url{https://cp-algorithms.com/num_methods/binary_search.html}
\item CSES 1085 \emph{Array Division} (practice) --- \ref{ex:arrdiv} yourself;
CSES 1620 \emph{Factory Machines} likewise. \newline\url{https://cses.fi/problemset/task/1085}
\item cppreference, ``std::lower\_bound'' (reference) --- the
library realization of \S\ref{sec:bs-template}. \newline\url{https://en.cppreference.com/w/cpp/algorithm/lower_bound}
\item Bloch, ``Nearly All Binary Searches and Mergesorts are Broken''
(blog, intermediate) ---
the midpoint-overflow bug this book's template avoids.
\newline\url{https://research.google/blog/} (post by Joshua Bloch, 2006)
\end{itemize}

\begin{rulebox}[Watch instead]
\smallskip\noindent\textbf{Best first.}
``Binary Search Algorithm'' --- Computerphile, 19 min ---
\url{https://youtu.be/hDn8iOc30Tk}. The halving
invariant of \S\ref{sec:bs-template} explained visually; Dr.\ Pound
indexes from 1 where the book uses 0 --- the boundary logic is
identical.
\smallskip\noindent
``Bug in Binary Search'' --- Computerphile, 12 min ---
\url{https://youtu.be/_eS-nNnkKfI}. The
\texttt{(lo+hi)/2} overflow of the Common mistakes section --- the real
Java library bug, worth watching before you write \texttt{mid} once.
Access date 2026-10-10.
\end{rulebox}
